\section{Nombres rationnels}
Les paramètres de procédures et les variables {\idt X} que nous avons utilisés jusqu'à 
présent ont des valeurs entières comprises entre $-32768$ et $32767$, ce que l'on indique 
par {\idt X $:$ integer} dans leur déclaration. On peut également utiliser des paramètres de 
procédures ou des variables ayant des valeurs entières comprises entre 
$-2 147 483 648$ et $2 147 483 647$, ce qu'il faut indiquer par {\idt X $:$ longint} dans la 
déclaration. Ces deux types d'entiers sont codés en base deux dans la mémoire de l'ordinateur, 
l'un sur 16 bits et l'autre sur 32.

Nous pouvons également utiliser en \pascal\ des nombres décimaux avec une virgule comme {\idt 3.14159} (de type {\idt real}). Ils
s'écrivent selon la convention anglo-saxonne  avec un point à la place de notre 
virgule. Les mathématiciens appellent ces nombres des {\em nombres rationnels}\index{Rationnel}. En \pascal, on
distingue les nombres entiers des nombres rationnels par la présence du point. Par exemple
1000 est un entier alors que $1000.0$ est un nombre rationnel. Dans l'écriture d'un nombre 
rationnel, il faut toujours écrire au moins
un chiffre avant et après le point. Par exemple $.0$ et $1.$ sont incorrects. Il faut écrire
$0.0$ et $1.0$.

Les nombres entiers et rationnels sont codés de manières différentes dans la mémoire 
de l'ordinateur. 

\section{Paramètres et variables de type ``real''}
On indique que les valeurs d'un paramètre formel ou d'une variable {\idt X}
sont des nombres rationnels par la déclaration {\idt X $:$ real}\footnote{\ Le choix du terme
``real'' (réel en français) est malheureux dans la mesure où il ne s'agit pas de nombres
réels puisque les nombres de type {\idt real} ne peuvent pas avoir une infinité de chiffres 
différents de zéro après la virgule$\:$!}.  Le paramètre effectif
correspondant peut être un nombre rationnel ou entier, le compilateur se chargeant dans ce
dernier cas de coder l'entier comme doivent l'être les nombres rationnels.

Les procédures de commande du \robot\ {\idt lg}, {\idt lgX}, {\idt lgY},  {\idt translater} et 
{\idt deplacer} ont des paramètres formels de type {\idt real} (et peuvent donc avoir des 
paramètres effectifs qui sont entiers ou rationnels).

\begin{exemple}[Spirale arithmétique]{\rm
La spirale ci-dessous est constituée d'arcs de cercles de rayons successifs
$100$$\:$; $100 / 1.5$ = $66.666\ldots$$\:$; $66.666\ldots / 1.5$ = $44.444\ldots$$\:$;
$44.444\ldots / 1.5$ = $29.629\ldots$$\:$; etc~:

 {\raggedright 
 \vbox to 0pt{\includevspacePostScript[4cm]{SpiraleArithmetique}{0.95}{0.95}}
 \input SpiraleArithmetique.tex
 }
\noindent $\Box$ }\end{exemple}

\section{Expressions rationnelles}\index{Expression rationnelle (réelle)}
L'addition ($+$), la soustraction ($-$) et la multiplication ($\ast$) de nombres rationnels se notent
comme pour les entiers. Par contre on utilise la barre de fraction $/$ pour la division de nombres 
rationnels. Par exemple $(((2 * (10.3 + 0.7)) - 1.0) / 3) = 7.0$. Si $r$ est un rationnel (ou un 
entier), {\idt sqr$(r)$} est son carré $r^{2}$ (de l'anglais \underline{sq}ua\underline{r}e, 
carré), {\idt sqrt$(r)$} est sa racine carrée $\sqrt{r}$ (de l'anglais \underline{sq}uare 
\underline{r}oo\underline{t}, racine carrée), {\idt abs$(r)$} est sa valeur \underline{abs}olue,  
{\idt trunc$(r)$} est sa partie entière (de l'anglais \underline{trunc}ate, tronquer) et 
{\idt round$(r)$} est l'entier le plus proche (de l'anglais \underline{round}, arrondir). Par exemple
{\idt trunc$(0.3) = $ trunc$(0.5) = $ trunc$(0.9) = 0$} tandis que
{\idt round$(0.3) = 0$} et {\idt round$(0.5) = $ round$(0.9) = 1$}.

Dans une expression rationnelle, on peut utiliser la longueur horizontale {\idt valLgX} et la
longueur verticale {\idt valLgY} d'un carreau de la grille de déplacement du \robot. Par exemple
{\idt lgX$($valLgX\ $\ast$\ 2$)$} double la taille horizontale d'un carreau de la grille.

\begin{exemple}[Spirale d'or]{\rm
La spirale d'or ci-dessous est constituée d'arcs de cercles successifs dont les rayons
ont pour rapport le {\em nombre d'or\/}
$\Phi = \frac{1}{2}(1 + \sqrt{5}) \simeq 1.618033989$. Pour les tenants d'esthétique
traditionnelle, la spirale d'or constituerait une figure géométrique très harmonieuse
à l'œil.

{\raggedright 
 \vbox to 0pt{\includevspacePostScript[4cm]{SpiraleDor}{0.95}{0.95}}
 \input SpiraleDor.tex
 }
\noindent $\Box$ }\end{exemple}

\section{Erreurs d'arrondi}
Pour les nombres réels\index{Réel}, qui ont un nombre infini de chiffres après la virgule, comme
$\frac{1}{3} = 0.3333333333333333\ldots$ ou $\pi = 3.1415926535897932\ldots$
l'ordinateur ne peut garder dans sa mémoire que quelques chiffres après la virgule, ce
qui conduit, comme pour une calculette, à des erreurs d'arrondi dans les calculs. Par
exemple en ma\-thé\-ma\-ti\-ques on a~:
\[\left(\left(\left(\left(\left(\sqrt{\sqrt{\sqrt{\sqrt{\sqrt{2}}}}}\right)^{2}\right)^{2}\right)^{2}\right)^{2}\right)^{2} = 
\left(\left(\left(\left(\sqrt{\sqrt{\sqrt{\sqrt{2}}}}\right)^{2}\right)^{2}\right)^{2}\right)^{2} =
\] \[\left(\left(\left(\sqrt{\sqrt{\sqrt{2}}}\right)^{2}\right)^{2}\right)^{2} =
\left(\left(\sqrt{\sqrt{2}}\right)^{2}\right)^{2} =
\left(\sqrt{2}\right)^{2} =
2\] 
tandis que sur
certains ordinateurs on peut trouver que~:

{\idt sqr$($sqr$($sqr$($sqr$($sqr$($sqrt$($sqrt$($sqrt$($sqrt$($sqrt$(2))))))))))$ =

$1.999999999999999997$}.

\noindent  On retrouve des problèmes d'arrondi similaires avec la grille de dé\-pla\-ce\-ment du \robot. De
manière idéale, les coordonnées cartésiennes du \robot\ devraient être réelles. 
L'ordinateur en conserve en mémoire une approximation rationnelle. Sur l'écran elles doivent
être arrondies à l'entier le plus proche pour correspondre à un pixel de l'écran.

\noindent \rule{137mm}{0.1mm} %---------------------------------------------
\begin{exercice}{Écrire des programmes \pascal\  pour réaliser les motifs de la figure
ci-dessous (seul le motif central qui utilise les fonc\-tions tri\-go\-no\-mé\-tri\-ques in\-tro\-dui\-tes
dans l'addenda du chapitre \ref{ch:rec}  est difficile).
\noindent $\Box$}\end{exercice}
\begin{figure}[htbp]
 \includevspacePostScript[4.6cm]{MotifsDeBeard}{0.95}{0.95}
 \caption{Motifs de Horemis et Beard$\:$; spirale de Baravelle}\label{fig:beard}
\end{figure}
\vfill
\begin{corrige}{
Nous donnons un programme, les autres se trouvant sur disquette~:
\input CarresImbriquesPeints.tex
\noindent $\Box$}\end{corrige}

%%%%%%%%%%%%%%%%%%%%%%%%%%%%%%%%%%%%%%%%%%%%%%%%%%%%%%%%%%%%%%%%%%%%%%%%%%%
\noindent \framebox[\textwidth]{
\begin{minipage}{\myparboxwidth}
\setlength{\parindent}{2em}
\small \sf
\vspace{\abovedisplayskip}
\section*{Addenda}
Les nombres rationnels sont qualifiés de réels en \pascal\ par abus de langage. Ils
peuvent se noter en utilisant la {\em notation scientifique}\index{Notation scientifique}, qui
consiste à faire
suivre la partie fractionnaire par un facteur d'échelle formé de la lettre {\idt e} ou
{\idt E} (qui se lit ``dix à la puissance'') suivie d'un exposant. 
  
\vspace{\abovedisplayskip}
\end{minipage}
}
%%%%%%%%%%%%%%%%%%%%%%%%%%%%%%%%%%%%%%%%%%%%%%%%%%%%%%%%%%%%%%%%%%%%%%%%%%%
%%%%%%%%%%%%%%%%%%%%%%%%%%%%%%%%%%%%%%%%%%%%%%%%%%%%%%%%%%%%%%%%%%%%%%%%%%%
\noindent \framebox[\textwidth]{
\begin{minipage}{\myparboxwidth}
\setlength{\parindent}{2em}
\small \sf
\vspace{\abovedisplayskip}
\noindent Par exemple $-3.0${\idt e}$2$ = $-300.0$,
$3.0${\idt e}$0$ =$3.0$ et $3.0${\idt e}$-2$ = $0.03$. Un nombre réel ne contient pas de blanc. Par conséquent
$-3.0\hspace{0.3ex}\blanc\hspace{0.3ex}${\idt e}$\hspace{0.3ex}\blanc\hspace{0.3ex}2$ est erroné. Sur les micros-ordinateurs, les réels sont 
généralement représentés avec 7 ou 8 chiffres significatifs en base 10, le plus petit réel 
strictement positif étant de l'ordre de $1.175${\idt e}$-38$ et le plus grand de l'ordre de
$3.403${\idt e}$38$.

\smallskip

\index{Expression rationnelle (réelle)}\index{Expression entière}La valeur d'une expression arithmétique 
peut être de type {\idt longint} ou {\idt real}, selon les règles suivantes~:
\begin{list}{$\scriptscriptstyle\bullet$}{
}
\item Si toutes les constantes et variables de l'expression sont entières (de type {\idt integer}
ou {\idt longint}) et si tous les
opérateurs de l'expression sont à résultats entiers (comme $+$, $-$,
$\ast$, {\mcl div}, {\mcl mod}, {\idt sqr}\footnote{\ Nous utilisons des typographies différentes
pour les {\em identificateurs réservés}\index{Identificateur réservé} (comme {\mcl program}
ou {\mcl mod})
et les {\em identificateurs prédéfinis}\index{Identificateur prédéfini}
(comme {\idt integer} ou {\idt sqr}). Les identificateurs introduits dans un programme \pascal\ 
(comme les noms de variable, de procédure, etc.) doivent être
différents des identificateurs réservés mais peuvent être identiques aux identificateurs prédéfinis. Cette
pratique est déconseillée car les programmes qui changent le sens habituel des identificateurs 
prédéfinis sont illisibles.}, {\idt abs} si leurs arguments sont entiers ou
{\idt trunc} et {\idt round}) alors l'expression est à résultat
entier de type {\idt longint}.
\item Si l'expression contient une constante réelle (avec un point ou un exposant {\idt e} ou
{\idt E}) ou une variables réelle (figurant dans une déclaration {\idt identificateur $:$ real})
ou un opérateur à résultat réel (comme $+$, $-$,
$\ast$, {\idt sqr} ou {\idt abs} si l'un des arguments est réel ou
 $/$ ou {\idt sqrt}) alors l'expression est à résultat
de type {\idt real}.
\end{list}

\smallskip

En \pascal, les paramètres effectifs doivent avoir le même type que les paramètres
formels, sauf dans les cas suivants~:
\begin{list}{$\scriptscriptstyle\bullet$}{
}
\item Si le type du paramètre formel est {\idt integer} alors le paramètre effectif peut
être de type {\idt longint} à condition que sa valeur soit comprise entre $-32768$ et $32767$$\:$;
\item Si le type du paramètre formel est {\idt real} alors le paramètre effectif peut
être de type {\idt longint} ou {\idt integer}.
\end{list}
 Dans les deux cas le compilateur se charge de coder l'entier 
 comme doivent l'être les grands entiers ou les nombres rationnels.
  
\vspace{\abovedisplayskip}
\end{minipage}
}
%%%%%%%%%%%%%%%%%%%%%%%%%%%%%%%%%%%%%%%%%%%%%%%%%%%%%%%%%%%%%%%%%%%%%%%%%%%
